% ===========================================================================
%  樊启斌《高等代数典型问题与方法》 · 第 4 章 矩阵 —— 秩论深化提取
%  （服务于考研蓝本第 4 章「矩阵的秩」，对应考纲〔矩阵〕〔向量〕〔线性方程组〕）
%
%  源：PDF p230–p247 = 印刷页 221–238（实测顶部页码：p230=221、p231=222、
%      p232=223、p235=226、p238=229、p241=232、p242=233、p243=234、
%      p245=236、p246=237、p247=238，即 PDF 页 = 印刷页 + 9，与 AGENTS.md §3.4 一致）
%
%  覆盖：§4.4 矩阵的秩（印刷页 221–236，基本理论 4.4.1–4.4.4 全收；
%        例题只取承载可复用结论/方法的，数字演算一律不进正文）
%        §4.5 矩阵的等价标准形（印刷页 236–238，基本理论 4.5.1 全收；
%        例 4.84–4.87 只取可复用的方法与结论）
%        §4.6 综合性问题（印刷页 238 起）：本片只到 p247（印刷 238），
%        该页仍属 §4.5 的例 4.86、4.87，故未涉及 §4.6 正文。
%
%  提取原则：只取定义 / 定理 / 性质 / 重要结论 / 方法要点 / 关键证明思路。
%            以下三块已由主控在 deep_ch04_authored.tex 写过，本片不重复：
%            「秩的四种等价定义」「秩-零化度定理」「秩不等式（含 Sylvester
%            推导、|r(A)-r(B)|≤r(A+B)）」「r(A^TA)=r(A) 与秩 1 矩阵完全刻画」。
%            伴随矩阵（含 r(A*) 三分式）见第 2 章，本片不重复。
% ===========================================================================


% ---------------------------------------------------------------------------
% 【深化】服务考点：秩的子式定义、可逆乘法下的秩不变性（考纲「矩阵」）
% 来源：樊 §4.4.1，印刷页 221（PDF p230）
% ---------------------------------------------------------------------------
\begin{fdeep}{子式怎么"数"秩：两条判别式 + 可逆乘法不变性}
$A$ 为 $m\times n$ 矩阵，任取 $k$ 行 $k$ 列（$k\le\min\{m,n\}$），交叉处的 $k^{2}$ 个元素按原次序构成的 $k$ 阶行列式，称为 $A$ 的 $k$ 阶子式。由此得到\textbf{两条可直接使用的判别式}：
\begin{itemize}[leftmargin=2.2em,itemsep=0.22em]
\item 若 $A$ 中\textbf{存在}一个 $r$ 阶子式不等于零，则 $r(A)\ge r$（给出秩的\textbf{下界}）；
\item 若 $A$ 中\textbf{所有} $r$ 阶子式全为零，则 $r(A)<r$（给出秩的\textbf{上界}）。
\end{itemize}
当 $r(A)=r$ 时，$r+1$ 阶子式（若存在）全为零。此外，$P,Q$ 分别为 $m$ 阶、$n$ 阶可逆矩阵时恒有
\fml{r(PA)=r(AQ)=r(A).}
\textbf{【反哺点】}服务考纲〔矩阵〕"理解矩阵的秩的概念"。考场上一旦要用"秩是几"去说理（而不是去算），就走子式这条路：要证 $r(A)\ge r$ 就\textbf{找出一个非零的 $r$ 阶子式}，要证 $r(A)<r$ 就\textbf{说明所有 $r$ 阶子式为零}。而 $r(PA)=r(AQ)=r(A)$ 是"乘可逆矩阵不改变秩"的完整写法，凡遇到可逆因子都可以放心消去。
\end{fdeep}


% ---------------------------------------------------------------------------
% 【深化】服务考点：分块矩阵的秩（考纲「矩阵」）
% 来源：樊 §4.4.3(4)–(7)，印刷页 221（PDF p230）
% ---------------------------------------------------------------------------
\begin{fdeep}{分块矩阵的秩：两条上界、一条等式、一条下界}
设 $A$ 为 $m\times n$ 矩阵，$B$ 为相应阶数的矩阵，则
\fml{\operatorname{rank}(A,\;B)\le \operatorname{rank}A+\operatorname{rank}B,}
\fml{\operatorname{rank}\begin{pmatrix}A\\ B\end{pmatrix}\le \operatorname{rank}A+\operatorname{rank}B;}
\textbf{而把两块放在对角线上时，秩是严格相加的}：
\fml{\operatorname{rank}\begin{pmatrix}A&O\\ O&B\end{pmatrix}=\operatorname{rank}A+\operatorname{rank}B;}
\textbf{放在三角位置的"下界"}（$C$ 与 $A$ 同宽）：
\fml{\operatorname{rank}\begin{pmatrix}A&O\\ C&B\end{pmatrix}\ge \operatorname{rank}A+\operatorname{rank}B.}
\textbf{【反哺点】}服务考纲〔矩阵〕"理解矩阵的秩的概念"。一句话记住：\textbf{横向拼、纵向叠都会"打折"（$\le$ 相加），只有对角摆放"不打折"（$=$ 相加）}；三角摆放至少不小于相加。做选择题估算秩的取值范围时，这四条就是边界。
\end{fdeep}


% ---------------------------------------------------------------------------
% 【深化】服务考点：求矩阵的秩的方法选择（考纲「矩阵」）
% 来源：樊 §4.4.4，印刷页 221（PDF p230）
% ---------------------------------------------------------------------------
\begin{fdeep}{求矩阵秩的五种方法：先看条件，再选工具}
书中列出的五条路子，实际对应五类题面条件：
\begin{enumerate}[leftmargin=2.4em,itemsep=0.22em]
\item \textbf{利用向量组的秩}——题面给的是"向量组""极大无关组""线性表示"；
\item \textbf{利用齐次方程组的理论}——题面给的是"$Ax=0$ 有非零解""解空间维数"；
\item \textbf{利用矩阵秩的等式或不等式}——题面给的是"$AB=O$""$A^{2}=A$"等代数关系；
\item \textbf{利用分块矩阵}——待证的式子是"几个秩的和差"，先把它们拼进一个分块矩阵；
\item \textbf{利用子空间的维数}，特别是线性变换\textbf{值域与核}的维数（详见第 5、6 章）。
\end{enumerate}
\textbf{【反哺点】}服务考纲〔矩阵〕"掌握用初等变换求矩阵的秩"。求"一个具体矩阵的秩"用化阶梯形；而遇到\textbf{含字母关系式的秩}（$AB=O$、$A^{2}=A$、$AB=BA$ 之类），初等变换帮不上忙，必须按上表选路子——绝大多数情况落在第 ③ 与第 ④ 条（见下面"分块初等变换"一块）。
\end{fdeep}


% ---------------------------------------------------------------------------
% 【深化】服务考点：秩不等式的取等条件（考纲「矩阵」）
% 来源：樊 例 4.58，印刷页 223（PDF p232）
% ---------------------------------------------------------------------------
\begin{fdeep}{秩不等式的取等条件：以 $r(A-B)\ge r(A)-r(B)$ 为例}
由 $A=(A-B)+B$ 及 $r(X+Y)\le r(X)+r(Y)$，立即得
\fml{\operatorname{rank}(A-B)\ge \operatorname{rank}A-\operatorname{rank}B.}
若 $A$ 可逆，则上面的等号成立\textbf{当且仅当} $BA^{-1}B=B$。
证明思路：作分块初等变换
\fml{\begin{pmatrix}A-B&O\\ O&B\end{pmatrix}\longrightarrow\begin{pmatrix}A&B\\ B&B\end{pmatrix}\xrightarrow{\ \text{行}\ }\begin{pmatrix}A&B\\ O&B-BA^{-1}B\end{pmatrix}\xrightarrow{\ \text{列}\ }\begin{pmatrix}A&O\\ O&B-BA^{-1}B\end{pmatrix},}
初等变换不改变秩，故 $r(A-B)+r(B)=r(A)+r(B-BA^{-1}B)$，等号成立等价于 $r(B-BA^{-1}B)=0$。
\textbf{【反哺点】}服务考纲〔矩阵〕"理解矩阵的秩的概念"。结论本身（下界式）已经够用，但\textbf{"取等条件"才是压轴题的考点}：题目问"何时 $r(A-B)=r(A)-r(B)$"，答案必须写成矩阵等式 $BA^{-1}B=B$，而不能只说"$B$ 的列都在 $A$ 的列空间里"。
\end{fdeep}


% ---------------------------------------------------------------------------
% 【深化】服务考点：证明秩等式/不等式的统一手法（考纲「矩阵」「线性方程组」）
% 来源：樊 例 4.58（p223）、例 4.60（p223–224）、例 4.66（p226–227）、
%        例 4.78 方法 3（p233）、例 4.79 方法 2（p234）、例 4.81（p234–235）、
%        例 4.82（p235）
% ---------------------------------------------------------------------------
\begin{fdeep}{分块初等变换：秩论证明的"统一手法"}
\textbf{做法}：把要证的几个秩所对应的矩阵块拼成一个 $2\times2$ 分块矩阵，然后\textbf{左乘或右乘可逆的分块矩阵}（即分块的行变换、列变换），秩始终不变；同一矩阵化成两种不同形状，就同时得到它的两种"秩的表达式"，比较即得等式或不等式。
\begin{itemize}[leftmargin=2.2em,itemsep=0.22em]
\item \textbf{得到等式}（打洞成对角块）：如上文 $r(A-B)+r(B)=r(A)+r(B-BA^{-1}B)$，把非对角块消成 $O$ 后，秩就"裂开"成两块之和。
\item \textbf{得到不等式}（Frobenius，例 4.79 方法 2）：
\fml{\begin{pmatrix}ABC&O\\ O&B\end{pmatrix}\xrightarrow{\ \text{行}\ }\begin{pmatrix}ABC&AB\\ O&B\end{pmatrix}\xrightarrow{\ \text{列}\ }\begin{pmatrix}O&AB\\ BC&B\end{pmatrix}}
（第一步：第一分块行加上"$A$ 乘第二分块行"；第二步：右乘可逆矩阵 $\begin{pmatrix}-E&O\\ C&E\end{pmatrix}$。）两端秩相等，而右端 $\ge r(AB)+r(BC)$，即得 $r(ABC)+r(B)\ge r(AB)+r(BC)$。
\end{itemize}
\textbf{【反哺点】}服务考纲〔矩阵〕"理解矩阵的秩的概念"+〔线性方程组〕。考研里凡出现"证明 $r(\cdot)\pm r(\cdot)\ge/\;=r(\cdot)$"的题，\textbf{第一反应就是分块初等变换}——它把"构造一个恒等式"变成了纯机械的消元操作，比凑向量组关系稳妥得多。
\end{fdeep}


% ---------------------------------------------------------------------------
% 【深化】服务考点：Sylvester 型秩恒等式（"打洞"）（考纲「矩阵」）
% 来源：樊 例 4.60 及其【注】，印刷页 223–224（PDF p232–233）
% ---------------------------------------------------------------------------
\begin{fdeep}{Sylvester 型秩恒等式：$A,D$ 可逆时的"打洞"公式}
设 $A$ 为 $m$ 阶、$D$ 为 $n$ 阶可逆矩阵，$B$ 为 $m\times n$、$C$ 为 $n\times m$ 矩阵，则
\fml{\operatorname{rank}A-\operatorname{rank}(A-BD^{-1}C)=\operatorname{rank}D-\operatorname{rank}(D-CA^{-1}B).}
证明只用可逆分块矩阵
$\begin{pmatrix}E&O\\ -CA^{-1}&E\end{pmatrix}$ 与 $\begin{pmatrix}E&-A^{-1}B\\ O&E\end{pmatrix}$ 把 $\begin{pmatrix}A&B\\ C&D\end{pmatrix}$ 分别化成 $\begin{pmatrix}A&O\\ O&D-CA^{-1}B\end{pmatrix}$ 与 $\begin{pmatrix}A-BD^{-1}C&O\\ O&D\end{pmatrix}$，比较两端即得。
\textbf{特例（书中【注】）}：取 $A=\lambda_0E_m$、$D=\lambda_0E_n$（$\lambda_0$ 为任意非零常数），则
\fml{m-\operatorname{rank}(\lambda_0E_m-BC)=n-\operatorname{rank}(\lambda_0E_n-CB).}
\textbf{【反哺点】}服务考纲〔矩阵〕"理解矩阵的秩的概念"、〔特征值〕。这条把 $BC$ 与 $CB$ 的秩"绑"在一起：$AB$ 与 $BA$ 的秩虽未必相等，但\textbf{它们"缺多少到满秩"是同一个数}。考研里"求 $r(E-AB)$"型的题、以及判断 $AB$ 与 $BA$ 特征值关系的题，都可以由这个特例一步跨过去。
\end{fdeep}


% ---------------------------------------------------------------------------
% 【深化】服务考点：用"方程组同解"证秩相等（考纲「线性方程组」「向量」）
% 来源：樊 例 4.62（方法 2），印刷页 224（PDF p233）
% ---------------------------------------------------------------------------
\begin{fdeep}{证秩相等的一条捷径：两个方程组同解}
设 $A$ 为 $m\times n$、$B$ 为 $n\times s$ 矩阵。若 $\operatorname{rank}(AB)=\operatorname{rank}B$，则对任一 $s\times t$ 矩阵 $C$，有
\fml{\operatorname{rank}(ABC)=\operatorname{rank}(BC).}
两个证法对应两种语言：\textbf{分块初等变换}（把 $\begin{pmatrix}ABC&O\\ O&B\end{pmatrix}$ 化两种形状）；\textbf{方程组同解}——设 $S_1,S_2$ 分别是 $ABCx=0$ 与 $BCx=0$ 的解空间，则 $S_1\supseteq S_2$；又任取 $x_0\in S_1$，有 $(AB)(Cx_0)=0$，由 $ABx=0$ 与 $Bx=0$ 同解得 $B(Cx_0)=0$，故 $x_0\in S_2$，于是 $S_1=S_2$，$\dim S_1=\dim S_2$，即得秩相等。
\textbf{【反哺点】}服务考纲〔线性方程组〕"理解同解方程组的概念"。\textbf{秩相等 $\Longleftrightarrow$ 解空间相同}是一枚硬币的两面：算秩算不动时，改证"两个方程组同解"，往往一行就完事；反过来题目说"同解"，也立刻知道秩相等、解空间维数相等。
\end{fdeep}


% ---------------------------------------------------------------------------
% 【深化】服务考点：矩阵方程有解的秩判据（考纲「线性方程组」「矩阵」）
% 来源：樊 例 4.71 及其【注】，印刷页 229（PDF p238）；§4.5.1(3)，印刷页 236（PDF p245）
% ---------------------------------------------------------------------------
\begin{fdeep}{矩阵方程有解：$r(A)=r(A,B)$，以及 $A=ABC$ 的充要条件}
\textbf{一、矩阵方程有解的判据}。设 $A$ 为 $m\times n$ 矩阵，$B$ 为 $m\times l$ 矩阵，$C$ 为 $p\times n$ 矩阵，则
\fml{AX=B\ \text{有解}\iff \operatorname{rank}A=\operatorname{rank}(A,B),}
\fml{XA=C\ \text{有解}\iff \operatorname{rank}A=\operatorname{rank}\begin{pmatrix}A\\ C\end{pmatrix}.}
\textbf{二、$A=ABC$ 的充要条件}。设 $A$ 为 $m\times n$、$B$ 为 $n\times m$ 矩阵，则存在 $m\times n$ 矩阵 $C$ 使 $A=ABC$，当且仅当
\fml{\operatorname{rank}A=\operatorname{rank}(AB).}
\textbf{必要性}只用到 $r(ABC)\le r(AB)$；\textbf{充分性}可以这样看：把 $A$ 按列分块 $A=(\alpha_1,\dots,\alpha_n)$，条件 $r(AB)=r(A)=r(AB,A)$ 说明方程组 $ABx=\alpha_j$ 对每个 $j$ 都有解（因增广矩阵 $\operatorname{rank}(AB,\alpha_j)\le\operatorname{rank}(AB,A)=r(AB)$），取这些解作列即得 $C$。
\textbf{【反哺点】}服务考纲〔线性方程组〕"理解非齐次线性方程组有解的判定条件"。第一条是\textbf{矩阵形式的"有解判别定理"}：$r(A)\ne r(A,B)$ 就是"系数矩阵秩小于增广矩阵秩"，比逐行消元快得多，选择题里几乎必用。
\end{fdeep}


% ---------------------------------------------------------------------------
% 【深化】服务考点：Frobenius 不等式及 $A+B$ 型秩不等式（考纲「矩阵」）
% 来源：樊 例 4.66，印刷页 226–227（PDF p235–236）；例 4.79，印刷页 233–234（PDF p242–243）；
%        例 4.81，印刷页 234–235（PDF p243–244）
% ---------------------------------------------------------------------------
\begin{fdeep}{Frobenius 不等式：秩不等式的"母版"}
设 $A$ 为 $m\times n$、$B$ 为 $n\times s$、$C$ 为 $s\times t$ 矩阵，则
\fml{\operatorname{rank}(ABC)\ge \operatorname{rank}(AB)+\operatorname{rank}(BC)-\operatorname{rank}B.}
它把 Sylvester 不等式 $r(AB)\ge r(A)+r(B)-n$ 推进到了\textbf{"三个因子"}的情形。证明有两条路：\textbf{满秩分解} $B=GH$（$G$ 列满秩、$H$ 行满秩）再套 Sylvester 不等式；或直接用分块初等变换。书中的【注】特别指出，\textbf{这条不等式的证明是各校考研的"高频"题}。
\textbf{【反哺点】}服务考纲〔矩阵〕"理解矩阵的秩的概念"。Frobenius 是\textbf{一切三因子秩不等式的出发点}：把 $ABC$ 中的某个因子取成 $O$，就退化成 $AB=O$、$BA=O$ 型结论；反复套用还能证明"秩序列 $r(A^{n})=r(A^{n+1})$"（见后文"秩序列的稳定性"一块）。记住它的\textbf{形状}——"中间项的秩要扣掉"——选择题里足以直接判对错。
\end{fdeep}


% ---------------------------------------------------------------------------
% 【深化】服务考点：$A+B$ 的秩的两个加强式（考纲「矩阵」）
% 来源：樊 例 4.66，印刷页 226–227（PDF p235–236）；例 4.81，印刷页 234–235（PDF p243–244）
% ---------------------------------------------------------------------------
\begin{fdeep}{$r(A+B)$ 的两个加强式：扣掉"重叠部分"}
一般情形只有 $r(A+B)\le r(A)+r(B)$。若 $A,B$ 之间有额外的交换性，上界可以\textbf{扣得更紧}（例 4.66，$AB=BA$）：
\fml{\operatorname{rank}(A+B)\le \operatorname{rank}A+\operatorname{rank}B-\operatorname{rank}(AB).}
另一个方向也有下界（例 4.81，$A,B$ 同形）：
\fml{\operatorname{rank}(A+B)\ge \operatorname{rank}(A,B)+\operatorname{rank}\begin{pmatrix}A\\ B\end{pmatrix}-\operatorname{rank}A-\operatorname{rank}B.}
后者由 Frobenius 不等式取 $P=(E_m,E_m)$、$M=\begin{pmatrix}A&O\\ O&B\end{pmatrix}$、$Q=\begin{pmatrix}E_n\\ E_n\end{pmatrix}$ 得到。
\textbf{【反哺点】}服务考纲〔矩阵〕"理解矩阵的秩的概念"。这两个式子说明：\textbf{$r(A+B)$ 的大小由 $A,B$ 的"重叠程度"决定，而 $r(AB)$ 与 $r(A,B)$ 正是重叠量的两种度量}——$AB$ 大（列空间重叠多）则 $A+B$ 的秩被压得低，$r(A,B)$ 大（列空间合起来张得宽）则 $A+B$ 的秩被抬得高。
\end{fdeep}


% ---------------------------------------------------------------------------
% 【深化】服务考点：秩序列一旦"停住"就永远停住（考纲「矩阵」「特征值」）
% 来源：樊 例 4.80，印刷页 234（PDF p243）；例 4.70(2)，印刷页 228（PDF p237）
% ---------------------------------------------------------------------------
\begin{fdeep}{秩序列的稳定性：$r(A^{k})=r(A^{k+1})$ 之后全部相等}
设 $A$ 为方阵，若存在正整数 $k$ 使
$\operatorname{rank}A^{k}=\operatorname{rank}A^{k+1}$，则对任意正整数 $n$ 有
\fml{\operatorname{rank}A^{k}=\operatorname{rank}A^{k+n}.}
特别地，对任意 $n$ 阶方阵 $A$，恒有 $\operatorname{rank}A^{n}=\operatorname{rank}A^{n+1}$。
证明可用 Frobenius 不等式对 $n$ 作归纳：$r(A^{k+n})=r(A^{n-1}A^{k}A)\ge r(A^{n-1}A^{k})+r(A^{k}A)-r(A^{k})=r(A^{k+n-1})+r(A^{k+1})-r(A^{k})$，结合归纳假设与题设即得下界，再由 $r(A^{k+n})\le r(A^{k})$ 收口。
另一条配套结论（例 4.70）：若 $A^{k-1}\alpha\ne0$ 而 $A^{k}\alpha=0$，则 $\alpha,A\alpha,\dots,A^{k-1}\alpha$ 线性无关；由此可证 $Ax=0$ 与 $A^{n+1}x=0$ 同解，从而 $r(A^{n})=r(A^{n+1})$。
\textbf{【反哺点】}服务考纲〔矩阵〕、〔特征值〕。"秩序列 $r(A),r(A^{2}),r(A^{3}),\dots$ 单调不增，且在 $n$ 步内稳定"是\textbf{幂零指数、最小多项式次数、Jordan 块结构}等一连串问题的统一入口；考场上遇到"$r(A^{k})=r(A^{k+1})$"就先把它翻译成"$A$ 的幂零部分已经全部消失"。
\end{fdeep}


% ---------------------------------------------------------------------------
% 【深化】服务考点：幂等矩阵的秩刻画（考纲「矩阵」「相似」）
% 来源：樊 例 4.83 及其【注】，印刷页 235–236（PDF p244–245）；
%        例 4.82，印刷页 235（PDF p244）
% ---------------------------------------------------------------------------
\begin{fdeep}{幂等矩阵：$A^{2}=A$ 的秩刻画与"秩可加"}
\textbf{一、最干净的等价刻画}（$A$ 为 $n$ 阶）：
\fml{A^{2}=A\iff n=\operatorname{rank}A+\operatorname{rank}(E-A).}
必要性由 $A(E-A)=O$ 得 $r(A)+r(E-A)\le n$，又 $n=r\bigl(A+(E-A)\bigr)\le r(A)+r(E-A)$；充分性可由两条不等式反推。
\textbf{二、相似标准形}：$A^{2}=A$ 且 $r(A)=r$ 时，存在 $n$ 阶可逆矩阵 $T$ 使
\fml{T^{-1}AT=\begin{pmatrix}E_r&O\\ O&O\end{pmatrix}.}
即\textbf{幂等矩阵一定可对角化}，且对角线上只有 $1$ 和 $0$，$1$ 的个数恰为 $r(A)$。
\textbf{三、两个配套结论}（例 4.82）：
\fml{\operatorname{rank}(A-ABA)=\operatorname{rank}A+\operatorname{rank}(E_n-BA)-n;}
配以 $A+B=E_n$ 与 $r(A)+r(B)=n$，可推出 $A^{2}=A$、$B^{2}=B$、$AB=O=BA$。更一般地，$A=\sum_{i}A_{i}$ 时
\fml{A^{2}=A\ \text{且}\ \operatorname{rank}A=\sum_{i}\operatorname{rank}A_{i}\iff A_{i}^{2}=A_{i}\ \text{且}\ A_{i}A_{j}=O\ (i\ne j).}
\textbf{【反哺点】}服务考纲〔矩阵〕"理解矩阵的秩"、〔特征值和特征向量〕"理解矩阵可相似对角化的充分必要条件"。幂等题在考场上只有两条路：\textbf{用 $n=r(A)+r(E-A)$ 算秩，或用 $\operatorname{diag}(E_r,O)$ 直接写出标准形}；最后一条"秩可加 $\iff$ 幂等且两两正交"是判断"能否拆成投影之和"的标准答案。
\end{fdeep}


% ---------------------------------------------------------------------------
% 【深化】服务考点：矩阵的等价标准形（考纲「矩阵」）
% 来源：樊 §4.5.1，印刷页 236（PDF p245）
% ---------------------------------------------------------------------------
\begin{fdeep}{矩阵的等价标准形：三件套}
\textbf{(1) 标准形}：设 $A$ 为 $m\times n$ 矩阵，$r(A)=r$，则存在 $m$ 阶可逆矩阵 $P$ 与 $n$ 阶可逆矩阵 $Q$，使
\fml{PAQ=\begin{pmatrix}E_r&O\\ O&O\end{pmatrix}.}
\textbf{(2) 满秩分解}：设 $A$ 为 $m\times n$ 矩阵，$r(A)=r$，则存在 $m\times r$ 的\textbf{列满秩}矩阵 $B$ 与 $r\times n$ 的\textbf{行满秩}矩阵 $C$，使
\fml{A=BC.}
\textbf{(3) 矩阵方程有解的判据}：$AX=B$ 有解 $\iff r(A)=r(A,B)$；$XA=C$ 有解 $\iff r(A)=r\begin{pmatrix}A\\ C\end{pmatrix}$（即把 $C$ 接在 $A$ 下方后秩不变）。
\textbf{【反哺点】}服务考纲〔矩阵〕"了解矩阵等价的概念"。这三条是\textbf{同一个事实的三种说法}：任何矩阵都"等价于"一个只由秩决定的标准形。(1) 用于\textbf{把一般矩阵换成对角块}再谈秩与构造；(2) 用于\textbf{把 $A$ 降维成"瘦长 $\times$ 扁平"}，是 $r(AB)\ge r(A)+r(B)-n$ 与 Frobenius 的证明引擎；(3) 是\textbf{矩阵形式的有解判别定理}，与第 5 章非齐次方程组判别式完全同源。
\end{fdeep}


% ---------------------------------------------------------------------------
% 【深化】服务考点：行满秩 / 列满秩矩阵（考纲「矩阵」「向量」）
% 来源：樊 例 4.85 及其【注 1】【注 2】，印刷页 237（PDF p246）；
%        例 4.85(1) 另见印刷页 237
% ---------------------------------------------------------------------------
\begin{fdeep}{行满秩与列满秩：单侧逆与消去律}
$A$ 为 $m\times n$ 矩阵，$r(A)=m$ 时称 $A$ \textbf{行满秩}（行向量组线性无关）；$r(A)=n$ 时称 $A$ \textbf{列满秩}。
\textbf{标准形}：$A$ 行满秩时，存在 $n$ 阶可逆矩阵 $Q$ 使
\fml{A=(E_m,\ O)Q;}
$A$ 列满秩时，存在 $m$ 阶可逆矩阵 $P$ 使 $A=P\begin{pmatrix}E_n\\ O\end{pmatrix}$。
\textbf{单侧逆}：$A$ 行满秩 $\iff$ 存在 $n\times m$ 矩阵 $B$ 使 $AB=E_m$（右逆）；$A$ 列满秩 $\iff$ 存在左逆。
\textbf{消去律}（书中【注 2】）：
\begin{itemize}[leftmargin=2.2em,itemsep=0.22em]
\item $A$ \textbf{行满秩}且 $GA=HA$，则 $G=H$（\textbf{右消去律}）；
\item $A$ \textbf{列满秩}且 $AG=AH$，则 $G=H$（\textbf{左消去律}）。
\end{itemize}
\textbf{【反哺点】}服务考纲〔向量〕"理解向量组线性相关、线性无关的概念"、〔矩阵〕。这条把"线性无关/相关"翻译成了\textbf{能不能约掉}：题目中若出现"由 $GA=HA$ 推出 $G=H$"，判分依据就是 $A$ 行满秩（列满秩取转置同理）。此外"行满秩 $\Rightarrow$ 有右逆"是解矩阵方程、构造反例时最常用的手段。
\end{fdeep}


% ---------------------------------------------------------------------------
% 【深化】服务考点：满秩分解的用法（考纲「矩阵」「线性方程组」）
% 来源：樊 例 4.87，印刷页 238（PDF p247）
% ---------------------------------------------------------------------------
\begin{fdeep}{满秩分解 $A=BC$ 的用法：$AX=0$ 与 $CX=0$ 同解}
称 $A=BC$ 为 $A$ 的\textbf{满秩分解表达式}，其中 $B$ 为 $m\times r$ 列满秩矩阵、$C$ 为 $r\times n$ 行满秩矩阵，$r=r(A)$。
\textbf{构造}：由 $A=P\begin{pmatrix}E_r&O\\ O&O\end{pmatrix}Q$ 直接拆开：
\fml{A=P\begin{pmatrix}E_r\\ O\end{pmatrix}\bigl(E_r,\ O\bigr)Q,}
取 $B=P\tbinom{E_r}{O}$（列满秩）、$C=(E_r,O)Q$（行满秩）即可。
\textbf{核心性质}：$A=BC$ 为满秩分解时，\textbf{$AX=0$ 与 $CX=0$ 是同解方程组}——因为 $B$ 列满秩 $\Rightarrow BX=0$ 只有零解，故 $BCx=0\iff Cx=0$。
\textbf{【反哺点】}服务考纲〔线性方程组〕"理解齐次线性方程组的基础解系"。要说清"$A$ 的零空间是谁"，\textbf{把 $A$ 写成 $BC$ 是最省力的一步}：$B$ 只负责"拉伸"，$C$ 决定了整个零空间。求基础解系时先做行变换取 $C$（阶梯形的非零行），本质上就是这条性质。
\end{fdeep}


% ---------------------------------------------------------------------------
% 【深化】服务考点：$r(E_m+AB)=r(E_n+BA)$（考纲「矩阵」）
% 来源：樊 例 4.77 及其【注】，印刷页 232（PDF p241）
% ---------------------------------------------------------------------------
\begin{fdeep}{$r(E_m+AB)=r(E_n+BA)$：加法扰动下的秩不变}
设 $A$ 为 $m\times n$、$B$ 为 $n\times m$ 矩阵，则\textbf{无需任何可逆性假设}，恒有
\fml{\operatorname{rank}(E_m+AB)=\operatorname{rank}(E_n+BA).}
证明用分块初等变换即可：
\fml{\begin{pmatrix}E&O\\ -B&E\end{pmatrix}\begin{pmatrix}E&-A\\ B&E\end{pmatrix}\begin{pmatrix}E&A\\ O&E\end{pmatrix}=\begin{pmatrix}E&O\\ O&E+BA\end{pmatrix},\qquad
\begin{pmatrix}E&A\\ O&E\end{pmatrix}\begin{pmatrix}E&-A\\ B&E\end{pmatrix}\begin{pmatrix}E&O\\ -B&E\end{pmatrix}=\begin{pmatrix}E+AB&O\\ O&E\end{pmatrix}.}
左端乘积都是可逆矩阵乘可逆矩阵，因而秩相等。
\textbf{【反哺点】}服务考纲〔矩阵〕"理解矩阵的秩"。这是"$AB$ 与 $BA$ 的秩可以不同，但\textbf{加上单位矩阵后秩就一样了}"的精确表述（与前面 Sylvester 型秩恒等式是同一件事的两种包装）。凡题目问"$r(E-AB)$ 与 $r(E-BA)$"的关系，答案永远是\textbf{相等}；若还给出两秩之和为 $n$，即可反推 $A$ 可逆。
\end{fdeep}


% ---------------------------------------------------------------------------
% 【深化】服务考点：互素多项式与秩的可加性（考纲「矩阵」「特征值」）
% 来源：樊 例 4.76，印刷页 231–232（PDF p240–241）
% ---------------------------------------------------------------------------
\begin{fdeep}{互素多项式与秩的相加：$f(A)=O$ 的分解}
设 $A$ 为 $n$ 阶方阵，$f(x)$ 为多项式且 $f(A)=O$。若 $f=f_1f_2$ 且 $(f_1,f_2)=1$（互素），则
\fml{\operatorname{rank}\bigl(f_1(A)\bigr)+\operatorname{rank}\bigl(f_2(A)\bigr)=n.}
\textbf{证明思路}：由互素，存在 $u(x),v(x)$ 使 $u f_1+v f_2=1$，代入 $A$ 得 $u(A)f_1(A)+v(A)f_2(A)=E$，于是 $r(f_1(A))+r(f_2(A))\ge n$；又 $f(A)=f_1(A)f_2(A)=O$ 给出 $r(f_1(A))+r(f_2(A))\le n$。两向夹逼即得等号。
\textbf{示例}：$A^{3}=E$ 时取 $f=x^{3}-1=(x-1)(x^{2}+x+1)$，两根互素，立即得
\fml{\operatorname{rank}(A-E)+\operatorname{rank}(A^{2}+A+E)=n.}
\textbf{【反哺点】}服务考纲〔矩阵〕、〔特征值〕。只要题目给出"$A$ 满足某个多项式方程 $f(A)=O$"，就可以把 $f$ \textbf{因式分解}，\textbf{互素的因子直接对应秩的相加}。这条与最小多项式、特征子空间直和分解是一条线上的事，是"由矩阵方程推秩等式"的通用机器。
\end{fdeep}


% ---------------------------------------------------------------------------
% 【深化】服务考点：实矩阵的共轭秩相等（考纲「特征值」）
% 来源：樊 例 4.75，印刷页 231（PDF p240）
% ---------------------------------------------------------------------------
\begin{fdeep}{实矩阵的共轭秩相等：$r(A-\mathrm{i}E)=r(A+\mathrm{i}E)$}
设 $A$ 为 $n$ 阶\textbf{实}方阵，$\mathrm{i}$ 为虚数单位，则
\fml{\operatorname{rank}(A-\mathrm{i}I)=\operatorname{rank}(A+\mathrm{i}I).}
证明只需构造一个同构：记 $U=\{x\mid (A-\mathrm{i}I)x=0\}$、$W=\{x\mid (A+\mathrm{i}I)x=0\}$。若 $\alpha+\mathrm{i}\beta\in U$（$\alpha,\beta$ 为实列向量），则 $(A-\mathrm{i}I)(\alpha+\mathrm{i}\beta)=0$，两边取共轭得 $(A+\mathrm{i}I)(\alpha-\mathrm{i}\beta)=0$，即 $\alpha-\mathrm{i}\beta\in W$。反之亦然，故映射
\fml{\varphi:\ U\longrightarrow W,\qquad \alpha+\mathrm{i}\beta\longmapsto\alpha-\mathrm{i}\beta}
是一一的线性映射（同构），于是 $\dim U=\dim W$，由解空间维数相同即得秩相等。
\textbf{【反哺点】}服务考纲〔特征值和特征向量〕"理解矩阵特征值和特征向量的概念"。考生常疑惑"实矩阵为什么会有复特征值、复特征向量成对出现"——答案就是\textbf{取共轭是同构}。由此 $A$ 的复特征值必共轭成对、实 Jordan 块成对出现，凡涉及"复特征值是否成对"的判断题都迎刃而解。
\end{fdeep}


% ---------------------------------------------------------------------------
% 【深化】服务考点：$A^{\mathrm T}B+B^{\mathrm T}A=O$ 时的秩下界（考纲「矩阵」「向量」）
% 来源：樊 例 4.68，印刷页 226–227（PDF p235–236）
% ---------------------------------------------------------------------------
\begin{fdeep}{实矩阵条件 $A^{\mathrm T}B+B^{\mathrm T}A=O$ 下的秩下界与取等}
设 $A,B$ 都是 $m\times n$ \textbf{实}矩阵，满足 $A^{\mathrm T}B+B^{\mathrm T}A=O$，则
\fml{\operatorname{rank}(A+B)\ge \max\{\operatorname{rank}A,\ \operatorname{rank}B\},}
且等号成立的\textbf{充分必要条件}是存在 $m$ 阶方阵 $P$，使 $B=PA$ 或 $A=PB$。
\textbf{证明思路（解空间的包含关系）}：设 $W_1,W_2,W_3$ 分别是 $Ax=0$、$Bx=0$、$(A+B)x=0$ 的解空间，则 $W_1\cap W_2\subseteq W_3$。由条件得 $(A+B)^{\mathrm T}(A+B)=A^{\mathrm T}A+B^{\mathrm T}B$，故对任意 $\xi\in W_3$ 有
\fml{0=\xi^{\mathrm T}(A+B)^{\mathrm T}(A+B)\xi=(A\xi)^{\mathrm T}(A\xi)+(B\xi)^{\mathrm T}(B\xi),}
于是 $A\xi=0$ 且 $B\xi=0$，即 $W_3\subseteq W_1\cap W_2$，从而 $W_3=W_1\cap W_2$。由 $\dim W_3\le\dim W_1,\dim W_2$ 即得下界；等号时 $W_3=W_1$，得 $W_1\subseteq W_2$，再由"解空间包含 $\iff$ 列空间包含"给出 $B=PA$。
\textbf{【反哺点】}服务考纲〔向量〕"理解向量组的秩与其行（列）向量组的秩之间的关系"、〔矩阵〕。这条把\textbf{秩的不等式换成了"解空间（列空间）的包含关系"}——凡是"实矩阵 $+$ 转置条件"的秩题，先把二次型 $(A+B)^{\mathrm T}(A+B)$ 写出来"配成平方和"，就能把包含关系两侧打通；这是用内积证秩不等式的标准动作。
\end{fdeep}


% ---------------------------------------------------------------------------
% 【深化】服务考点：满秩的稳定性、由行（列）和条件直接读出秩（考纲「矩阵」）
% 来源：樊 例 4.55，印刷页 222（PDF p231）；例 4.56，印刷页 222–223（PDF p231–232）；
%        例 4.64(2)，印刷页 226（PDF p235）
% ---------------------------------------------------------------------------
\begin{fdeep}{满秩是"稳定"的：多项式扰动下秩不会掉}
$A,B$ 为 $n$ 阶复方阵，$B$ 满秩，则存在正实数 $a$，当 $0<s<a$ 时 $C=A+sB$ 仍为满秩矩阵。
\textbf{方法}：由 $B$ 满秩得 $C=(AB^{-1}+sE)B$，于是 $|C|=|AB^{-1}+sE|\cdot|B|$。而 $f(s)=|AB^{-1}+sE|$ 是 $s$ 的 $n$ 次多项式，\textbf{多项式的根至多 $n$ 个}，故在 $s=0$ 附近（足够小的正数区间内）$f(s)\ne0$，即 $C$ 满秩。
\textbf{【反哺点】}服务考纲〔矩阵〕"理解矩阵的秩"。这是一条"用多项式的根有限"证明\textbf{满秩在扰动下保持}的典范：凡是题目把可逆矩阵"动一下"（$A+\varepsilon B$、$A+sB$、$A+tE$），都可以先把行列式看成某个变元的多项式，再用"根是有限个"一句话定性。它也解释了为什么\textbf{可逆是"开条件"、不可逆是"闭条件"}。
\end{fdeep}


% ---------------------------------------------------------------------------
% 【深化】服务考点：行（列）和条件直接给出秩（考纲「矩阵」「线性方程组」）
% 来源：樊 例 4.56，印刷页 222–223（PDF p231–232）；例 4.64(2)，印刷页 226（PDF p235）
% ---------------------------------------------------------------------------
\begin{fdeep}{列和为零、行和为常数：直接读出秩的两种题面}
\textbf{一、列和全为零 $\Rightarrow$ 秩缺 $1$}（例 4.56）：若 $n$ 阶方阵 $A$ 满足 $a_{ii}>0$、$a_{ij}<0\ (i\ne j)$，且
\fml{\sum_{i=1}^{n}a_{ij}=0\qquad (j=1,2,\cdots,n),}
则 $\operatorname{rank}A=n-1$。
\textbf{方法}：列和为 $0$ 说明 $A$ 的列向量线性相关，故 $r(A)\le n-1$；再取 $a_{11}$ 的余子式 $M_{11}$，由
$a_{jj}=-\sum_{i\ne j}a_{ij}=\sum_{i\ne j}|a_{ij}|>\sum_{i=2,\,i\ne j}^{n}|a_{ij}|$
知删去第 1 行第 1 列后的 $n-1$ 阶子式\textbf{严格对角占优}，故 $M_{11}>0$，于是 $r(A)\ge n-1$。
\textbf{二、行和为 $1$ $\Rightarrow$ $(E-A)x=0$ 有非零解}（例 4.64(2)）：若 $\sum_{j=1}^{n}a_{ij}=1$，取 $\alpha=(1,1,\cdots,1)^{\mathrm T}$，则 $A\alpha=\alpha$，故 $(E-A)\alpha=0$，即 $r(E-A)<n$。
\textbf{【反哺点】}服务考纲〔矩阵〕、〔线性方程组〕。"行和/列和"是\textbf{免费送来的特征向量}：列和为 $0$ 立刻给出 $r(A)\le n-1$（并说明 $0$ 是特征值），行和为 $1$ 立刻给出 $r(E-A)<n$。读题时见到"每行元素之和为 $c$"，应当条件反射地写出 $A(1,\dots,1)^{\mathrm T}=c(1,\dots,1)^{\mathrm T}$，比消元快得多。
\end{fdeep}


% ---------------------------------------------------------------------------
% 【深化】服务考点：Sylvester 不等式的取等条件（考纲「矩阵」）
% 来源：樊 例 4.78 及其【注】，印刷页 233（PDF p242）
% ---------------------------------------------------------------------------
\begin{fdeep}{Sylvester 不等式取等的充要条件}
$A$ 为 $m\times n$、$B$ 为 $n\times s$ 矩阵时，$r(AB)\ge r(A)+r(B)-n$（推导见主控深化框）。进一步，\textbf{等号成立}
\fml{\operatorname{rank}(AB)=\operatorname{rank}A+\operatorname{rank}B-n}
\textbf{的充分必要条件是}：矩阵方程
\fml{AX+YB=E_n}
有解。证明思路是\textbf{等价标准形}：取可逆 $P,Q$ 使 $PAQ=\begin{pmatrix}E_r&O\\ O&O\end{pmatrix}$（$r=r(A)$），令 $Q^{-1}B=\begin{pmatrix}C\\ D\end{pmatrix}$，则 $PAB=\begin{pmatrix}C\\ O\end{pmatrix}$，于是 $r(AB)=r(C)$，而 $r(AB)\ge r(A)+r(B)-n$ 等价于"$C$ 的行数 $r$ 与 $B$ 的行秩 $r(B)$ 之差不超过 $D$ 的行数 $n-r$"，取等即 $C$ 的行线性无关、$D$ 的每一行都可由 $C$ 的行线性表示——这正是 $AX+YB=E_n$ 有解的等价说法。
\textbf{【反哺点】}服务考纲〔矩阵〕、〔线性方程组〕。知道"取等 $\iff$ $AX+YB=E_n$ 有解"后，判断"$r(AB)$ 能否取到最小值 $r(A)+r(B)-n$"就有了\textbf{可操作的判据}（转化为解一个矩阵方程），而不用再去凑基础解系。
\end{fdeep}


% ===========================================================================
% 以下按深化提取原则跳过（页码为印刷页）
% ===========================================================================
% fan p222 例 4.53（东南大学 2003，由 $PQ=O$ 的秩反推 $P$ 的秩）：结论就是
%          $PQ=O\Rightarrow r(P)+r(Q)\le n$ 的直接应用，已并入秩不等式一块，不另立。
% fan p222 例 4.54（北京科技大学 2008，$A^{2}=A,B^{2}=B,E-A-B$ 可逆 $\Rightarrow r(A)=r(B)$）：
%          技巧核心"右乘可逆矩阵秩不变"已见于"子式"与"可逆乘法不变性"一块，
%          结论只对该题成立，不另立块。
% fan p224 例 4.61（河南师范大学 2012，证明该伴随矩阵 $A^{*}$ 的秩为 $n$ 或 $1$）：
%          伴随矩阵的一般秩公式 $r(A^{*})=n/1/0$ 已收入第 2 章（deep_ch02_authored.tex），
%          本片不重复。
% fan p225 例 4.63（南京大学 2007，$r(V^{\mathrm T}A^{-1}U+E_m)<m\Rightarrow r(A+UV^{\mathrm T})<n$）：
%          其可复用内核是行列式恒等式 $|A+UV^{\mathrm T}|=|A|\cdot|E_m+V^{\mathrm T}A^{-1}U|$，
%          属行列式技巧；且"秩 $<n\iff$ 行列式为 $0$"已是常见结论，不另立块。
% fan p226 例 4.65（上海大学 2000，由 $AB$ 反求 $BA$）：$3\times2$ 与 $2\times3$ 的具体
%          数字演算，结论不可复用，跳过。
% fan p227 例 4.67（$AA^{\mathrm T}=a^{2}E\Rightarrow r(aE_n-A)=r\bigl((aE_n-A)^{2}\bigr)$）：
%          结论针对特殊矩阵，与考纲内的通用秩方法无接口，跳过。
% fan p228 例 4.69（华中师范大学 2008，$AB=O\Rightarrow r(A)\le[n/2]$，并构造达到 $[n/2]$ 的
%          例子）：$[n/2]$ 这一取整结论属竞赛型构造题，超出考纲要求的秩方法，跳过。
% fan p229 例 4.72（山东大学 2017，$\operatorname{tr}(A^{*})=a\ne0$ 反推 $r(A)=n-1$ 并求
%          $|\lambda E-A^{*}|$）：是 $r(A^{*})$ 三分式与降阶公式的应用，二者均已在前文/第 2 章，
%          不另立块。
% fan p229–230 例 4.73（华中师范大学 1999，$AC+BD=O$ 时 $r\tbinom{C}{D}=t\iff r(C)=t$）：
%          用到"满列秩矩阵可写成 $T\tbinom{E_s}{O}$"，该手法已并入"行满秩/列满秩"一块；
%          结论本身偏专，不另立块。
% fan p230 例 4.74（华中科技大学 2017、清华大学 2008，$AB=BA=O$ 且 $r(A^{2})=r(A)$
%          $\Rightarrow r(A+B)=r(A)+r(B)$）：方法完全被"等价标准形 + 分块初等变换"两块覆盖，
%          不另立块。
% fan p238 例 4.86（山东大学/中国科学技术大学，列满秩实矩阵补成可逆矩阵）：结论是
%          "列满秩可补全为基"的技术性引理，考纲内不直接考，跳过。
% fan p238 §4.6 综合性问题：本片范围到 p247（印刷 238）为止，该页仍是 §4.5 的例 4.86、
%          4.87，未进入 §4.6 正文，故本节无内容需要判断。
% ===========================================================================


% ===========================================================================
%  【主控裁决 · 第 4 章深化料取舍】（依据 AGENTS.md §2 决定二）
%  手写块 deep_ch04_authored.tex 已覆盖：秩的四种等价定义、秩-零化度定理、
%  秩不等式（含 Sylvester 推导）、r(A^TA)=r(A) 与秩 1 矩阵完全刻画。
%
%  【保留】（服务考纲"秩"与"秩与行列向量组秩的关系"）
%   樊 块3  求矩阵秩的五种方法：先看条件，再选工具   ← 方法选择判据，价值高
%   樊 块5  分块初等变换：秩论证明的统一手法
%   樊 块7  证秩相等的一条捷径：两方程组同解
%   樊 块8  矩阵方程有解 r(A)=r(A,B)；A=ABC 充要条件
%   樊 块9  Frobenius 不等式（秩不等式的母版）
%   樊 块13 矩阵的等价标准形：三件套
%   樊 块14 行满秩与列满秩：单侧逆与消去律
%   樊 块15 满秩分解 A=BC 的用法
%   樊 块22 Sylvester 不等式取等的充要条件
%   北大 块⑥ 行秩=列秩（子式证法五步）★考纲明确要求
%   北大 块⑦ 初等变换下的秩、阶梯形与极大无关组
%   北大 块③ 极大线性无关组与向量组的秩
%
%  【删除及理由】
%   · 北大 块⑧ 方阵情形 |A|=0 与齐次方程组 → 与第 3 章「可逆等价刻画」重复
%   · 北大 块① 线性相关性方程组判定 / 块② 替换定理 → 与 9 讲第 6 讲正文重复
%   · 北大 块④ 线性表出与方程组同解 → 与樊 块7 重复（取樊块更简洁）
%   · 北大 块⑤ 秩的子式定义怎么用 → 与手写块「四种等价定义」重叠
%   · 樊 块1  子式怎么数秩 → 同上，重叠
%   · 樊 块2  分块矩阵的秩 → 与第 3 章分块内容重复
%   · 樊 块4  秩不等式取等条件 → 与手写块「秩的不等式推导」重叠
%   · 樊 块6  Sylvester 型秩恒等式(打洞) → 偏技巧，考纲无接口
%   · 樊 块10 r(A+B) 的两个加强式 → 与手写块重叠
%   · 樊 块11 秩序列稳定性 r(A^k)=r(A^{k+1}) → 超纲（考纲未涉及幂序列）
%   · 樊 块12 幂等矩阵 A^2=A 的秩刻画 → 边缘，非考纲主线
%   · 樊 块16 r(E_m+AB)=r(E_n+BA) → 竞赛级
%   · 樊 块17 互素多项式与秩 → 超出考纲（属多项式理论接口）
%   · 樊 块18 共轭秩相等 → 复数域内容，考研不考
%   · 樊 块19 A^T B+B^T A=O 下的秩下界 → 竞赛级
%   · 樊 块20 满秩的稳定性(多项式扰动) → 超纲
%   · 樊 块21 列和为零/行和为常数读秩 → 与 9 讲正文重复
%
%  保留 12 块（樊 9 + 北大 3）。来源 36.8KB + 16.8KB → 精选后约 9KB。
% ===========================================================================
